\documentclass[11pt]{article}
\usepackage[T1]{fontenc}
\usepackage{lmodern,amsmath,amssymb,amsthm}
\usepackage[margin=1in]{geometry}
\usepackage[hidelinks]{hyperref}
\newtheorem{theorem}{Theorem}[section]
\newtheorem{proposition}[theorem]{Proposition}
\newcommand{\Jac}{\operatorname{Jac}}
\newcommand{\Pic}{\operatorname{Pic}}
\newcommand{\CH}{\operatorname{CH}}
\newcommand{\Q}{\mathbb Q}
\newcommand{\F}{\mathbb F}
\newcommand{\Z}{\mathbb Z}
\title{Modified diagonals and branch-divisor points on dihedral genus-two quotients}
\author{\small\href{mailto:admin@cosmosgeek.org}{admin@cosmosgeek.org}\qquad
\href{mailto:b321014@iiit-bh.ac.in}{b321014@iiit-bh.ac.in}}
\date{Working paper for specialist review}
\begin{document}\maketitle
\begin{abstract}
For a family of plane quintics with a dihedral action of order ten, we compute
the image of the canonical modified diagonal under a reflection correspondence.
Its pushforward to the genus-two quotient is an explicit multiple of a
branch-divisor class. This class has infinite order on an infinite sequence
of rational fibres. We determine the degree-two elliptic quotient locus
in a two-parameter chart. One component gives an explicit cover of every
prescribed non-CM elliptic curve after an extension of degree at most eight.
We find a nontorsion specialization and a quadratic character that can force
the trace of the projected point to vanish. A second family gives a
nontorsion cycle image on a geometrically simple Jacobian. In that family
we classify the degree-two and degree-three loci and give a local obstruction
to elliptic quotients with prescribed Tate parameter. For a fixed dihedral
quotient the moving branch class is computed by a real-multiplication isogeny.
\end{abstract}
\section{The reflection correspondence}
We use smooth projective curves in characteristic zero and Chow groups with
rational coefficients. Let $\pi:C\to D$ be the quotient by an involution
$\tau$, with ramification divisor $R$, branch divisor $B=\pi_*R$ and
$d=\deg R$. Let $e=K_C/(2g(C)-2)$ for $g(C)>1$. The modified diagonal
$\Delta_e\in\CH^2(C^3)$ is the small diagonal, minus the three partial
diagonals with $e$ in the remaining coordinate, plus the three copies of
$C$ with $e$ in each of the other two coordinates \cite{GS}. The insertions
are exterior products of cycles. Put
\[
 \mathcal T_\tau(\Delta_e)=\pi_*\operatorname{pr}_{3,*}
       (\operatorname{pr}_{12}^*\Gamma_\tau\cdot\Delta_e)
 \in\Pic^0(D)\otimes\Q .
\]
\begin{proposition}\label{graph}
The graph contraction is
\[
 \mathcal T_\tau(\Delta_e)=
 \left(1-\frac d{2g(C)-2}\right)B-
 \frac{2d}{2g(C)-2}K_D .
\]
\end{proposition}
\begin{proof}
Before the final pushforward, the seven signed terms are
$R,-de,-\tau_*e,-\tau_*e,e,e,0$. For the small diagonal, the graph of
$\tau$ meets the diagonal transversely at fixed points since the multiplier
there is $-1$. The final intersection vanishes in $\CH^3(C^2)$. Since
$\tau_*e=e$, the sum is $R-de$. Apply Hurwitz and
$\pi_*\pi^*K_D=2K_D$.
\end{proof}
Graph contractions also occur in Chow--Heegner point constructions
\cite{CH}. An explicit nontorsion image detects nonvanishing of the
modified diagonal modulo rational equivalence.

\section{Plane quintics of genus six}
Consider the smooth members of
\begin{equation}\label{quintic}
 C:\ z^5+w^5+2a_2z^2w^2t+2a_1zwt^3+2a_0t^5=0,\qquad a_0\ne0 .
\end{equation}
The involution $\tau$ exchanges $z$ and $w$. The rotation is
$(z,w,t)\mapsto(\zeta_5z,\zeta_5^{-1}w,t)$ over a field containing
$\zeta_5$. The coordinate vertices lie outside $C$, so the rotation acts
freely. The rotation quotient has genus two, with affine equation
$v^2=(a_2x^2+a_1x+a_0)^2-x^5$ for $x=zw,\ v=(z^5-w^5)/2$ on $t=1$.
Dihedral covers of this type occur in \cite{AP}.

The reflection quotient $D$ has invariant generators $u=z+w,x=zw$
on $t=1$. The equation for $x$ is
\[
 (5u+2a_2)x^2+(2a_1-5u^3)x+u^5+2a_0=0.
\]
With $Y=2(5u+2a_2)x+2a_1-5u^3$ this gives
\begin{equation}\label{sextic}
 Y^2=5u^6-8a_2u^5-20a_1u^3-40a_0u+4a_1^2-16a_0a_2 .
\end{equation}
Conversely $z,w$ are the roots of $Z^2-uZ+x$. The fixed locus lies on
$z=w$ together with $P=[1:-1:0]$. Along the fixed line its intersection
with $C$ is transverse since the tangent of $C$ at a fixed point has
multiplier $-1$. Thus $R\sim H+P$ for a line section $H$, and $g(D)=2$.
As $K_C=2H$, Hurwitz gives $\pi^*K_D\sim H-P$. With $Q=\pi(P)$ we get
\begin{equation}\label{class}
 B\sim2K_D+2Q,\qquad \delta=B-3K_D=2Q-K_D,\qquad
 \mathcal T_\tau(\Delta_e)=\tfrac25\delta .
\end{equation}
Here $u_Q=-2a_2/5$, $Y_Q=-2a_1-8a_2^3/25$: the affine product
$zw/t^2$ has a pole at $P$, selecting the infinite-$x$ sheet over
$5u+2a_2=0$.

\begin{theorem}\label{points}
For $a_2=0$, $a_1=s$, $a_0=1$ and any $s\in21\mathbb Z\setminus\{0\}$,
$\delta$ has infinite order in $\Jac(D_s)(\Q)$. The generic section over
$\Q(s)$ also has infinite order.
\end{theorem}
\begin{proof}
Here $D_s:Y^2=5u^6-20su^3-40u+4s^2$ and $Q_s=(0,-2s)$.
At 3 and 7 the plane quintic has smooth diagonal reduction, and the
sextic reduces to $5u(u^5-8)$. Its derivatives at zero and a nonzero
root are $-40$ and $200$, so the sextic has good reduction at both primes.
The hyperelliptic involution $\iota$ gives $\delta=Q_s-\iota Q_s$,
nonzero since distinct degree-one divisors cannot be equivalent on
a curve of positive genus. At both primes $Q_s$ reduces to the
Weierstrass point $(0,0)$, so $\delta$ reduces to zero. Prime-to-$p$
torsion injects under good reduction in residue characteristic $p$.
An element in both torsion kernels would therefore vanish. Specializing at
one such smooth fibre also rules out a finite annihilator of the generic
section.
\end{proof}

\section{The elliptic quotient locus}
The parameters $(a_2,a_1,a_0)$ have weights $(1,3,5)$. Set $a_2=1$,
$a=a_1$, $b=a_0$. Put
\begin{align*}
 D_0={}&108a^5+16a^4-900a^3b-128a^2b+2000ab^2+256b^2+3125b^3,\\
 F_1={}&16(a^2-4b)^2-3125b^3,\\
 F_2={}&-128(a^2-4b)^2+34000ab^2-37800a^3b+9261a^5-25000b^3 .
\end{align*}
\begin{proposition}\label{locus}
On $bD_0\ne0$ the degree-two elliptic quotient locus of \eqref{sextic}
over an algebraic closure is $F_1F_2=0$. Its two components are
geometrically irreducible rational curves.
\end{proposition}
\begin{proof}
The discriminant of $(x^2+ax+b)^2-x^5$ is $b^5D_0$ by the resultant.
An affine singularity of the quintic would give a repeated root of
that polynomial, by multiplying the first derivative equations by
$z,w$ respectively. The infinity points are smooth. For the binary
sextic $f$ in \eqref{sextic} define the covariants
\[
 I=(f,f)_4/86400,\quad L=(I,f)_4/2880,\quad
 M=(I,L)_2/24,\quad N=(I,M)_2/24.
\]
The unnormalized transvection $(g,h)_j$ is
$\sum_{i=0}^j(-1)^i\binom ji
(\partial_u^{j-i}\partial_v^i g)(\partial_u^i\partial_v^{j-i}h)$,
where $v$ homogenizes $u$. Differentiation in coefficient order
$u^2,uv,v^2$ gives $\det(L,M,N)=2^{17}D_0F_1F_2$.
This nonzero degree-15 invariant is proportional to $R$ of
\cite[Lemma 2.5]{MS}, by the one-dimensionality of the degree-15
invariant space \cite{CFG}. The criterion there gives precisely the locus
of the extra involution.

Writing $x=b/a^2$ and dividing by $a^4$ gives
\[
 16(1-4x)^2-3125a^2x^3,\qquad
 -25000x^3a^2+(34000x^2-37800x+9261)a-128(1-4x)^2 .
\]
The first quadratic has no root over $\mathbb C(x)$ by valuation at zero.
The second has nonsquare discriminant
$(20x-49)^2(20x-9)^2(441-1280x)$. Both are primitive over
$\mathbb C[x]$, and $a=0$ is not a component. Substitution gives dominant
rational parametrizations
\[
 F_1:\ h=4(v^2-20)/625,\qquad a=hv,\qquad b=5h^2.
\]
\[
 F_2:\ a=\frac{256(w+11)^2}{25(21-w)^3},\
 b=\frac{256(w+11)^4(w+21)}{3125(21-w)^5}.
\]
\end{proof}

The moduli image of \eqref{quintic} is dense in the Humbert surface
$H_5$. The endomorphism $S$ constructed below supplies the inclusion.
For a numerical certificate of the dimension use the absolute
invariants in the convention of \cite{Shaska}. At $(a,b)=(1,1)$
their values and derivatives from the sextic transvectants are
\[
 (i_1,i_2)=(1089/1369,27008289/1266325),
\]
\[
 \frac{\partial(i_1,i_2)}{\partial(a,b)}=
 \begin{pmatrix}
 -118800/50653&59400/50653\\
 1213071552/46854025&-284909724/46854025
 \end{pmatrix}.
\]
The determinant is $-41307216192/2565726409$. For reproduction,
the invariants $J_2=800(40b-3a^2)$, $J_4=40000(20b-9a^2)^2$ and
\[
 \frac{J_6}{320000}=600000b^3+18432ab^2+225000ab^3-466000a^2b^2
                -9216a^3b+119700a^4b+1152a^5-6399a^6
\]
give $i_1=144J_4/J_2^2$,
$i_2=-1728(J_2J_4-3J_6)/J_2^3$. The point is in $bD_0\ne0$.
Irreducibility of $H_5$ and this full-rank differential give the
assertion.

\subsection{A different multiplicative degeneration}
For the subfamily in Theorem~\ref{points},
\[
 \operatorname{disc}(5u^6-20su^3-40u+4s^2)
 =2^{18}5^5(108s^5+3125)^2.
\]
For $p\ge7$ a negative valuation of $s$ gives potentially good
reduction: adjoin $\alpha^3=s$ and put $u=\alpha x$, $Y=sy$. The
reduction $y^2=5x^6-20x^3+4$ is smooth. A multiplicative quotient
therefore requires a specialization near $s_*$ with
$108s_*^5+3125=0$. Put $A=-6s_*^2/25$, so $A^5=-2/3$,
$s_*=(5/2)A^3$. The limiting sextic is
\[
 5(u^2-Au-A^2)^2(u^2+2Au+5A^2).
\]
Let $\phi_\pm=(1\pm\sqrt5)/2$. At the two nodes $u=A\phi_i$,
the linear terms of the critical values in $s-s_*$ are
$-40A^3\phi_i(s-s_*)$, both units times $s-s_*$. The normalization
is the conic parametrized by
$u=A(t-t^{-1}-1)$, $V=A\sqrt5(t+t^{-1})$, where
$V=Y/(u^2-Au-A^2)$. The pairs above the nodes are
$(\phi_i+1\pm b_i)/2$, $b_i^2=3(\phi_i+2)$.
The graph is a bouquet with two edges of length $k=v(s-s_*)$ and
monodromy $kI_2$. In the split marking the off-diagonal period has
constant term the cross-ratio of these two pairs,
\[
 c=(-7+3\sqrt5)/2=-\phi_+^{-4}
\]
or its inverse for a change of orientation. Indeed choose
$b_+b_-=3\sqrt5$; the numerator and denominator cross-difference
products are $(-5+3\sqrt5)/2$ and $(-5-3\sqrt5)/2$.
Their product is a unit at the retained primes. The universal
unit argument given below for the three-edge graph applies here
with zero edge exponents. It yields $Q_{12}=c(1+O(s-s_*))$
with integral Taylor coefficients over the unramified splitting field.

On this lattice the RM endomorphism $S$ is symmetric for $kI_2$,
has trace $-1$, determinant $-1$, and is
$\left(\begin{smallmatrix}0&1\\1&-1\end{smallmatrix}\right)$
after signed permutation. Multiplicative adjunction gives
$Q_{11}/Q_{22}=Q_{12}$. Write $U=Q_{22}$, $V=Q_{12}$.
An elliptic quotient character $(m,n)$ sends the two period
columns to $U^mV^{m+n}$ and $U^nV^m$. As $v(U)>0$ and $v(V)=0$,
eliminating the target Tate period gives $V^{n^2+mn-m^2}=1$.
The exponent is nonzero for every nonzero integral $(m,n)$.
Conversely a torsion $V$ gives an elliptic splitting by raising both
coordinates to its order.

Write $s_p=v_p(F_{p-(5/p)})$ for the indicated Fibonacci number, with
$v_p(p)=1$. Binet's formula gives $v_p(\log c)=s_p$: with
$h=p-(5/p)$ even, $\phi_+^{2h}-1=\sqrt5\,\phi_+^h F_h$, and $h$ is a
$p$-adic unit. Thus $v_p(s-s_*)>s_p$ makes $\log Q_{12}$ nonzero and
the Jacobian geometrically simple. This test is specific to the
two-loop degeneration just computed.

\subsection{A specialization with two nontorsion projections}
Take $(a_2,a_1,a_0)=(1,0,256/3125)$ on $F_1$. The base quintic is
squarefree: a repeated nonzero root of $(x^2+a_0)^2-x^5$ would force
$x=16/25$ and $a_0=-256/3125$. Over $\Q(i,\sqrt5)$ set $u=2X/5$,
$Y=(8\sqrt5/125)y$, $X=2iz$, $y=8iv$. Equation \eqref{sextic} becomes
$v^2=z^3P_3(z+z^{-1})$, where $P_3(S)=S^3+2iS^2-3S-4i$. For
$\epsilon=\pm1$ the formulas
\[
 S=z+z^{-1},\qquad W_\epsilon=v(z+\epsilon)/z^2,\qquad
 W_\epsilon^2=(S+2\epsilon)P_3(S)
\]
give complementary elliptic quotients with deck involutions given by
reciprocal inversion and its composite with $\iota$. Set
$A_\epsilon=-2\epsilon+4i$, $x=A_\epsilon/(S+2\epsilon)$ and
$y=A_\epsilon W_\epsilon/(S+2\epsilon)^2$. The elliptic equation is
\[
 y^2=x^3+(9-8\epsilon i)x^2+(4-28\epsilon i)x-12-16\epsilon i .
\]
The marked point maps first to $(z,v)=(i/2,i\sqrt5/8)$, then to
\[
 P_\epsilon=\left((-8+4\epsilon i)/5,\ \sqrt5(-6\epsilon+8i)/25\right).
\]
The points at infinity on $D$ represent $K_D$ and have opposite images
$(0,\pm A_\epsilon)$. Thus $\delta$ maps to $2P_\epsilon$. At a place
over $2+i$, normalized by $v(5)=1$, the coordinates of $P_+$ are
$-4/(2+i)$ and $-2\sqrt5/(2+i)^2$. In the integral formal group
at infinity $T=-x/y$ has valuation $1/2$. For $v(T)>1/4$ the formal
logarithm has the same valuation as $T$: the coefficient of degree $n$
has valuation at least $-v_5(n)$ by integration of the invariant
differential. Hence $P_+$ has infinite order. Complex conjugation
takes $P_+$ to $-P_-$.

\subsection{A map to a prescribed elliptic curve}
\begin{theorem}\label{universal}
For every non-CM elliptic curve $E/K$ over a number field there is a
smooth member of \eqref{quintic} over an extension of degree at most two
whose reflection quotient admits a degree-two map to $E$ over an
extension of degree at most eight over $K$. Both elliptic projections
of $\delta$ have infinite order on the generic member of this
one-parameter construction.
\end{theorem}
\begin{proof}
In weighted parameters take
\[
 a_2=\frac{5(21-w)}4,\quad a_1=20(w+11)^2,\quad
 a_0=\frac{(w+11)^4(w+21)}4.
\]
The weight-15 homogenization of $D_0$ evaluates to
$(390625/32)(w+11)^{10}(w^2+135)^2$. Thus the base quintic is squarefree
for $w\ne-11,-21$ and $w^2+135\ne0$, including $w=21$.
Use $s=\sqrt5$ and substitute
\[
 u=(w+11)v,\quad d(z)=(s-1)z-(s+1),\quad
 v=-2(z+1)/d(z),\quad y=\frac{Yd(z)^3}{20(w+11)^2}.
\]
The resulting equation is $y^2=\kappa(Az^6+Bz^4+Cz^2+D)$ with
$\kappa=10(w+11)$ and
\begin{align*}
 A&=20(s-2),& B&=(5-s)w+15(s-1),\\
 D&=-20(s+2),& C&=(5+s)w-15(s+1).
\end{align*}
The quotient coordinates are $(z^2,y)$ and $(z^{-2},y/z^3)$,
with cubic coefficients $\kappa(A,B,C,D)$ and $\kappa(D,C,B,A)$.
The identities
\[
 B^2-3AC=10(3-s)(w^2+135),\qquad
 \operatorname{disc}(AX^3+BX^2+CX+D)=400(w^2+135)^2
\]
and their conjugates give $j(E_1)=j(E_2)=64(w^2+135)/5$. Choose
$w^2=5j(E)/64-135$ with sign avoiding $-11,-21$. Since $j(E)\ne0,1728$,
the quotient $E_1$ becomes isomorphic to $E$ after at most one quadratic
extension of $K(w,s)$.

At $w=11$ the marked point has $u_Q=-5$, $Y_Q=-19985$. Its transformed
coordinates are
\[
 z_Q=-\frac{49+5s}{49-5s},\qquad y_Q=\frac{7034720s}{(49-5s)^3}.
\]
On the monic integral model of $E_1$ the coordinates are
$x=\kappa A z_Q^2,\ y_m=\kappa A y_Q$. At the prime over $569$
dividing $49-5s$, the factors $\kappa A$ and the displayed numerators
are units. Since the norm of $49-5s$ is $4\cdot569$, one has
$v(x)=-2$, $v(y_m)=-3$. The formal logarithm applies at $v(-x/y_m)=1$.
The conjugation $s\mapsto-s$ exchanges the quotient equations and takes
$(z_Q,y_Q)$ to $(1/z_Q,y_Q/z_Q^3)$, proving that both quotient points
have infinite order. On either quotient $\iota$ acts as inversion, so
$\delta$ maps to twice the image of $Q$. Specialization proves the
generic assertion.
\end{proof}

\begin{proposition}\label{twist}
With $T=w^2+135$, put $E_w:y^2=x^3-3Tx-2wT$. Choose $r^2=\Delta(E)$
and $w=c_6(E)s/(8r)$ with the sign chosen as above. Over $K(w,s)$,
$E_1$ is the twist of $E$ by the squareclass
$d=25(w+11)(1-s)c_4(E)/r$. If $5$ is nonsquare in $K(w)$, every
$K(w,s)$-point of $E_1$ transported by the twisting isomorphism has
zero trace from $K(w,s,\sqrt d)$ to $K$.
\end{proposition}
\begin{proof}
The monic form for $E_1$ has coefficients
$a_2=\kappa B$, $a_4=\kappa^2AC$, $a_6=\kappa^3 A^2D$.
Its $c_4,c_6$ are those of $E_w$ multiplied by $d_1^2,d_1^3$ with
$d_1=\kappa(s-5)/3$. Also $T=5c_4(E)^3/(64\Delta(E))$, and the
usual invariant formulas give the twist of $E$ defining $E_w$ as
$3sc_4(E)/(2r)$. Their product is $d$. Write $F=K(w)$.
The nontrivial automorphism of $F(s)/F$ negates $s,r$ and gives
\[
 \operatorname{N}_{F(s)/F}(d)=2500(w+11)^2c_4(E)^2/\Delta(E).
\]
Its squareclass is $5$ in $F^\times/F^{\times2}$, since
$\Delta(E)=5(c_6(E)/(8w))^2$. Thus $d$ is nonsquare, and the nontrivial
quadratic automorphism acting on the transported point gives its inverse.
\end{proof}
For example, for $E:y^2+xy=x^3-x^2-8x-7$ over $K=\Q(\sqrt{-2})$,
the invariants are $393,7803,-109$. Thus $K(w)=\Q(\sqrt{-2},\sqrt{-545})$
with quadratic squareclasses $-2,-545,1090$. Proposition~\ref{twist}
applies. Matching the geometric elliptic quotient in this example
does not furnish a point of infinite order on the prescribed curve
over $K$.

\section{A genus-four cover and a simple specialization}
Consider $C_t:y^5=x^3(x-1)^2(x-t)$ with $t\ne0,1$. This cyclic cover
has four totally ramified points and genus four. The genus-four family
of \cite[Example 3.8]{LS} is obtained after changes of coordinate and
scalar: choosing $\eta^2=t(t-1)$ and $\beta^5=-t^2/\eta^3$, set
$z=(x-t-\eta)/(x-t+\eta)$, $w=\beta y/x$ and
$\rho=(t+\eta)/(t-\eta)$. Then
$w^5=((z+1)/(z-1))((z+\rho)/(z-\rho))^2$.
Adjoin $c$ with $c^5=t^4(t-1)^3$. A reflection is
\[
 \tau(x,y)=\left(\frac{t(x-1)}{x-t},
                    \frac{cx(x-1)}{(x-t)y}\right).
\]
It has two fixed points: over each of the two fixed points of the
base involution, inversion on the five cyclic sheets fixes one.
These base points are disjoint from the branch locus. Hence the
reflection quotient has genus two and $\deg B=2$.
Proposition~\ref{graph} gives
$\mathcal T_\tau(\Delta_e)=\frac23(B-K_D)$.

\subsection{The quotient model and its branch}
Put $a=-t(t-1)/c$, $b=-t^2(t-1)/c^2$. The canonical coordinates
$z_1,z_2,z_3,z_4$ given respectively by
$dx/y,\ x\,dx/y^2,\ x(x-1)\,dx/y^3,\ x^2(x-1)\,dx/y^4$ satisfy
\[
 z_1z_4=z_2z_3,\qquad
 z_1z_2^2-z_1^2z_3-(1-t)z_2z_4^2-tz_3^2z_4=0.
\]
Their regularity follows by checking orders at the four ramification
points and their independence by the four different rotation characters.
The ratios $z_2/z_1=x/y$ and $z_3z_1/z_2^2=(x-1)/x$ recover $x,y$.
Under $\tau$ the pairs $(z_1,az_4)$ and $(z_2,bz_3)$ are interchanged
up to the same projective scalar. Set
\[
 U=\frac{z_2+bz_3}{z_1+az_4},\qquad
 V=\frac{z_1-az_4}{z_1+az_4},\qquad
 W=\frac{z_2-bz_3}{z_1+az_4},\qquad T_0=VW.
\]
For a direct elimination put
\begin{align*}
 h&=b^2-aU,& P_0&=b^2+aU-2abU^2,& R_0&=2bU-b^2U^2-aU^3,\\
 A_0&=h^2-4ab(bU+1)^2,& B_0&=2(bU+1)(bP_0+aR_0).
\end{align*}
The two equations yield
\[
 V^2=\frac{P_0-2a(bU+1)T_0}{h},\quad
 W^2=\frac{R_0-2b(bU+1)T_0}{h},\quad A_0T_0^2+B_0T_0-P_0R_0=0.
\]
Thus $U,T_0$ generate the invariant field. With
$Y=(A_0T_0+B_0/2)/(bh)$, the last discriminant gives
$Y^2=a^2U^6+(4a^2/b)U^5-10aU^3+4bU+1$.
Adjoin $s$ with $s^3=a$, and put $X=sU$, $q=b/s$.
The normalized equation is then
\begin{equation}\label{four}
 D_q:\quad Y^2=X^6+(4/q)X^5-10X^3+4qX+1,\qquad q^3=t/(t-1).
\end{equation}
The identification just exhibited is over the extended parameter field.
The equation itself is defined over $\Q(q)$.

For generic $q$ the two branch points satisfy
\begin{align}
 &(1-4q^3)X^2-10q^2X+q^4-4q=0,\label{branchx}\\
 &Y=-\frac{(qX+1)(q^3+3qX-3q^2X^2-X^3)}{q(q^2-X)}.\notag
\end{align}
Indeed, at the two zeros of $A_0$ the pole sheet of $T_0$ gives
$Y=-B_0/(2bh)$. There $V^2$ generically has odd-order pole, so
the quadratic extension defined by $V$ ramifies there. Its total
ramification degree is two. Substituting $X=sU$ gives
\eqref{branchx}. In particular their sum defines $B$ over $\Q(q)$.
For $q=2$, the $X$-polynomial is $X^2+(40/31)X-8/31$ and the two
$Y$-values are given by $(13748X-15689)/8649$. These formulas give
the divisor with the indicated sign in the smooth fibre.

\begin{proposition}\label{fourpt}
At $q=2$, $B-K_D$ has infinite order over $\Q$ and $\Jac(D_2)$ is
geometrically simple. In particular the generic section has infinite
order over $\Q(q)$.
\end{proposition}
\begin{proof}
The two $X$-coordinates are distinct, so their sum cannot be a
hyperelliptic fibre. This proves $B-K_D\ne0$. The sextic is squarefree
in characteristics 3 and 11. At 3 the two points reduce to the two
points over $X=1$ with $Y=\pm\sqrt2$ (substitute
$X=(-20\pm18\sqrt2)/31$ and $Y=(-84591\pm27496\sqrt2)/29791$ with correlated signs before reducing).
This is a canonical divisor. Thus torsion in this class would be
3-primary. At 11 the counts of $D_2$ over $\F_{11}$ and $\F_{121}$ are
$8,148$. They give $\#\Jac(D_2)(\F_{11})=95$, prime to 3,
which proves the first assertion. For reproduction, the two counts are
$2+\sum_x\#\{y:y^2=x^6+2x^5-10x^3+8x+1\}$ in the respective fields;
one can use $\F_{11}[v]/(v^2-2)$ in the second sum.

The Frobenius polynomial is
$Z^4-4Z^3+21Z^2-44Z+121$. Its quadratic factors over $\Q(\sqrt5)$
have traces $2\pm\sqrt5$ and constant term 11. Their discriminants
$-35\pm4\sqrt5$ are totally negative, and a root $\pi$ recovers the
quadratic field through $\pi+11/\pi$. Hence the polynomial is
irreducible over $\Q$, and $11\nmid21$ gives ordinarity.
In the notation of \cite[Theorem 6]{HZ}, $a=-4,\ b=21$, so
$a\ne0$ and $a^2$ is different from $b+11,2b,3b-33$. The reduction is
absolutely simple. Specialization of geometric endomorphisms at good
reduction excludes a geometric elliptic factor in characteristic zero.
\end{proof}

The simple fibre has a consequence for target selection. For a fixed
integer $N$, only finitely many geometric elliptic $j$-invariants occur
as targets of degree $\le N$ from smooth members of \eqref{four}. For
each optimal degree $m\ge2$ the locus is the Humbert surface of
discriminant $m^2$ \cite{Kumar}. Its inverse image is a proper closed
subset of the parameter curve by Proposition~\ref{fourpt}. On each
of the finitely many fibres, the endomorphisms
$\phi^*\phi_*$ of its Jacobian for optimal covers $\phi$ of degree
$m\le N$ have Rosati squared norm $2m^2$ (use the trace on $H^1$).
A positive-definite lattice contains finitely many elements of bounded
norm, and their images specify the embedded elliptic factors. Nonoptimal
maps add isogenies of bounded degree and hence only finitely many
further targets. The explicit family in Theorem~\ref{universal} has
different moduli and is not subject to this finiteness conclusion.

\section{Small-degree quotients of the genus-four family's sextic}
Put $s=q^3$ and $w=s+s^{-1}$ in \eqref{four}. In the invariant
normalization of \cite{Shaska} its Igusa--Clebsch invariants are
\[
 J_2=1000,\quad J_4=62500,\quad
 J_6=18505000-1440000w,\quad J_{10}=12800000(w-2)^2 .
\]
These are obtained by the normalized transvections
$A=[f,f]_6,\ I=[f,f]_4,\ B=[I,I]_4,\ C=[I,[I,I]_2]_4$; here
$[g,h]_j$ is $(g,h)_j$ multiplied by
$(\deg g-j)!(\deg h-j)!/(\deg g!\deg h!)$. The formulas are
$J_2=-120A,\ J_4=-720A^2+6750B,\
J_6=8640A^3-108000AB+202500C$, and $J_{10}=\operatorname{disc}(f)$.
Thus the two equations of the moduli curve in the $J_2\ne0$ chart are
\[
 C_1=J_2^2-16J_4=0,\qquad C_2=(64J_6-J_2^3)^2-663552J_2J_{10}=0.
\]

\begin{proposition}
The geometric degree-two locus in the smooth family \eqref{four}
is $(s+1)(32s^2+3061s+32)=0$. Its geometric degree-three locus
has ten distinct $w$-values. Only $w=-2$ is rational. At the nine
other degree-three values the class $B-K_D$ is nontorsion.
\end{proposition}
\begin{proof}
With the covariants of Proposition~\ref{locus} the determinant becomes
$-64q^{-9}(s-1)^3(s+1)(32s^2+3061s+32)$. The same invariant criterion
proves the degree-two statement and $q=-1$ is the only rational $q$ in it.

For degree three we give elimination details to specify the parameter
charts. Use the normal forms in \cite{Shaska}, with absolute invariants
in that paper's convention. In the degenerate chart
$(3X^2+4)(X^3+X+c)$ the invariants are
\[
 J_2=774,\quad J_4=36(268+837c^2),\quad
 J_6=36(76760+290574c^2-729c^4),\quad J_{10}=432(4+27c^2)^3.
\]
The first two absolute invariants force $c^2=12353/13392$ and
$w=-111600508651/19559994112$. The third differs by
$537328796797921779/1530373478645898674176$.
In the generic chart take
\[
 f_1=X^3+aX^2+bX+1,\qquad f_2=4X^3+b^2X^2+2bX+1,\quad
 u=ab,\ v=b^3,\ H=v-9-2u .
\]
At $b=0$, $C_1=0$ gives $a^3=-2565/704$, where
$C_2=136650973826879014545/14641$. At $H=0,\ b\ne0$,
$C_1=12(v-27)^3(v^2-11v-32)/v$.
After dividing $v^4 C_2$ by $64(v-27)^8$, its remainder modulo
$v^2-11v-32$ is $166212(2864+1043v)$, which is coprime to that
quadratic. $v=27$ is singular.

On the remaining chart put $t=-H/(v-27)$ and set
\begin{align*}
 R&=4u^3-u^2v-18uv+4v^2+27v,&
 \eta&=-48R/(v(v-27)^3),\\
 D_t&=1+2t-t^2/2,& E_t&=t^2-4t-2,\\
 K_t&=(2t+1)(1+2t-2t^2),&
 Q_t&=32t^5+96t^4+56t^3+76t^2+22t+1 .
\end{align*}
Transvection of $f_1f_2$ gives the absolute invariants
\[
 i_1=\frac{9((t+1)^4+(t+4)\eta)}{4D_t^2},\qquad
 i_2=\frac{27((t+1)^6+\tfrac32(t-8)(t+1)^2\eta+\tfrac38\eta^2)}{8D_t^3},
 \quad i_3=-\frac{9\eta^3}{131072D_t^5}.
\]
For \eqref{four} these are
$9,-(37719+23328w)/3125,243(w-2)^2/39062500$, respectively.
The first equality gives $\eta=3K_t/(t+4)$; $t=-4$ fails and
$K_t=0$ would give a singular curve. The last two equations give
\[
 w(t)=2+\frac{3125(2t+1)Q_t}{256(t+4)^2E_t^3},\qquad
 Q_t^2-64(2t+1)(1+2t-2t^2)^3(t+4)E_t=0.
\]
The polynomial factors as
\[
 (4t+3)(4t^2-6t+3)(8t^3+36t^2+20t+3)
 (16t^4-32t^3+96t^2+82t+19).
\]
Every root is realized: solve
\[
 K_t v^2+(2t+1)(4t^3+78t^2-18t-11)v-8(t+4)(3t+1)^3=0
\]
and set $u=((t+1)v-27t-9)/2$, $b^3=v$, $a=u/b$.
The polynomial at $v=27$ equals $4(t+4)(9t+5)^2$.
The ten roots above avoid $t=0,-4,-1/2,-1/3,-5/9$ and
$K_tE_t=0$. Substitution in the quadratic gives
$R=(v-27)^2((t+1)^2(2t+1)v-2(3t+1)^3)/4\ne0$.
Thus the reconstruction gives a smooth normal form, and its
invariants determine the same geometric genus-two curve.
On $Y^2=f_1f_2$ the actual degree-three map is
\[
 U=X^2/f_1,\quad V=Y(X^3-bX-2)/f_1^2,\qquad
 V^2=-(R/v)U^3+2(6a^2-ab^2-9b)U^2+(b^2-12a)U+4.
\]

The first factor gives $w=-2$. For the remaining three irreducible
$t$-factors, elimination of $t$ against $w(t)$ gives
\begin{align*}
 W_2&=2476099w^2-4557521w+8976271,\\
 W_3&=20511149w^3-43927820019w^2+147887021288w-150744429817,\\
 W_4&=146211169851w^4-4170608746308w^3+50055219167049w^2\\
    &\hspace{22mm}-229010043991482w+384192134208241 .
\end{align*}
The first two are irreducible modulo 2, the third modulo 11. The two
distinct primes $3,7$ divide $W_2(-2)$, the primes $3,47$ divide
$W_3(-2)$ and $13,37$ divide $W_4(-2)$. All corresponding leading
coefficients are units. At primes of $\Q(w)$ over $(p,w+2)$, and then
of a field containing $q$, one has $s\equiv-1$. A pole-free expression
for the branches of \eqref{branchx}, putting $A=1-4s$, is
\begin{align*}
 X_\pm&=(5q^2\pm2(s+1)\sqrt q)/A,\\
 Y_\pm&=(164s^3+27s^2-138s-1\mp2q(4s^3-153s^2+107s+14)\sqrt q)/A^3 .
\end{align*}
At the specified primes their sum reduces to a canonical fibre: the
two coordinates are $X=q^2,\ Y=\pm4q\sqrt q$. The reduced sextic is
smooth by the discriminant formula. In characteristic zero the two
$X$-coordinates are distinct. As in Theorem~\ref{points}, two
different good residue characteristics rule out torsion. The
complementary elliptic projections give an isogeny, so at least one
projection is nontorsion at every one of these nine parameters.
\end{proof}

At the additional degree-two values the analogous nonvanishing holds
on both factors. To see the geometry, choose
\[
 s=(-3061+225\sqrt{185})/64,\quad
 z=-(8s+17)/225,\quad r=(5\sqrt{37}-13\sqrt5)/8.
\]
Under $X=(z+ru)/q,\ v=sY/r^3$, the reciprocal sextic has the form
$v^2=u^6+Au^5+Bu^4+Cu^3+Bu^2+Au+1$, where
$A=4\sqrt5+2\sqrt{37}$, $B=-23-2\sqrt{185}$,
$C=14\sqrt5+6\sqrt{37}$. The quotients are
$W_\pm^2=(U\pm2)(U^3+AU^2+(B-3)U+C-2A)$, with $U=u+u^{-1}$.
Direct quartic invariants give $j=4096$ for both. At primes of residue
characteristics 3 and 37, $s\equiv-1$ and the preceding specialization
argument applies. The involution negating $\sqrt5,\sqrt{37}$ fixes $s$
and exchanges the two quotients by $u\mapsto-u$, so both images are
nontorsion. All these observations concern geometric elliptic factors.

The finite calculation also determines the fields of the elliptic
invariants in the degree-three case. In the nondegenerate chart used
above set
\[
 T=6u^2/v-u-9,\quad W=1-12u/v,\quad R_0=R/v,\quad
 Z=4T^2+3R_0W .
\]
The cubic in $U,V$ gives $j_1=16v^2 Z^3/(R_0^2(v-27)^3)$.
With $z=2t$, the following coefficient data give a short certificate
for the two targets. Write $S_j=j_1+j_2$, $P_j=j_1j_2$.
Eliminating $v$ by the displayed quadratic gives
\[
\begin{array}{c|l}
 \text{minimal polynomial of }z&\text{reduced expressions}\\ \hline
 z^2-3z+3&S_j=251498-143759z,\quad P_j=-68433479-43892640z\\
 z^3+9z^2+10z+3&S_j=-259117-899290z-782739z^2\\
   &P_j=26283214+75435525z+53768088z^2\\
 z^4-4z^3+24z^2+41z+19 &j_1=j_2=486-1458z+1053z^2+189z^3 .
\end{array}
\]
These data can be reproduced by replacing $v^2$ by the linear polynomial
in $v$ from the reconstruction quadratic and adding or multiplying its
two conjugates. All are algebraic integers. No $j_i$ here is rational:
in the quadratic case the $z$ coefficient of $j^2-S_jj+P_j$
would force $j=3990240/13069$, which fails the constant equation.
In the cubic case the nonconstant coefficients require two different
rational values. In the quartic case the displayed representative has
nonconstant degree below four. At the remaining rational parameter
$q=-1$ a reciprocal substitution gives elliptic quartics
\[
 W_\pm^2=(S\pm2)(S^3-4S^2-3S-2)
\]
with invariants $j=-25/2,-121945/32$. All the elliptic factors
arising at these degree-two or degree-three fibres are therefore
potentially good at every $p\ge7$. Composing with further isogenies
preserves potential good reduction.


\section{A local obstruction to a prescribed multiplicative factor}
Let $p\ge7$, $L=\Q_p(\sqrt5)$ with its unramified valuation $v(p)=1$,
and retain the genus-two family \eqref{four}.
\begin{theorem}\label{local}
Let $E/\Q_p$ have multiplicative reduction, let $q_E$ be its Tate
parameter after the possible unramified splitting twist, and write
$e=v(q_E)$. If $p\ne67$ and $p\nmid e$, then no smooth geometric
member of \eqref{four} over $\overline{\Q}_p$ has a nonzero homomorphism
from its Jacobian to $E$. The family nevertheless has infinitely many
algebraic parameters with multiplicative elliptic quotients.
\end{theorem}

\subsection{The normalized multiplicative periods}
For $q$ a unit with $q^3-1$ a unit the discriminant in the preceding
section gives good reduction. If $v(q)>0$, set $\alpha^5=q$, $X=\alpha u$.
The reduced binary branch form is $Z(4u^5+Z^5)$ with distinct roots,
including the one at infinity. This gives potentially good reduction.
The reciprocal change $X\mapsto1/X,\ Y\mapsto Y/X^3$ replaces $q$ by
$1/q$. So a multiplicative quotient is possible only at $q^3\equiv1$.
After a cube-root-of-unity scaling take $q=1+h$, $v(h)>0$.

Put $g=(x-1)(x^2+3x+1)$ and label its roots
$a_0=1,\ a_1=(-3+\sqrt5)/2,\ a_2=(-3-\sqrt5)/2$. At $h=0$ the
components $y=\pm g$ meet at these three nodes. If $c_a(h)$ is the
critical point of $f_{1+h}$ near $a$, then
\[
 f_{1+h}(c_1(h))=(36/25)h^2+O(h^3),\qquad
 f_{1+h}(c_a(h))=4(a-a^5)h+O(h^2)\quad(a=a_1,a_2).
\]
These follow by expanding $g^2+4h(x-x^5/(1+h))$ and its first
derivative. Since $\operatorname N(a-a^5)=-45$, the two coefficients
are units for $p\ge7$. The edge lengths are $2v(h),v(h),v(h)$.
In the graph basis $e_0-e_1,e_0-e_2$ the valuation matrix is
$v(h)M_0$ with $M_0=\left(\begin{smallmatrix}3&2\\2&3\end{smallmatrix}\right)$.

Denote the symmetric, principally polarized multiplicative period matrix
by $Q=(Q_{ij})$ in this basis. We record the computation of its first
units, since valuations alone would miss the obstruction. Taking branch
coordinates with first jets $x-a$ on both components gives plumbing
parameters $y_a=-f_{1+h}(c_a(h))/(4g'(a)^2)+\cdots$. Thus
$y_{a_0}\sim-9h^2/625$ and $y_a\sim3(7a+3)h/25$ for $a=a_1,a_2$.
The stable forms on the plus component are
$d\log((x-a_0)/(x-a_i))$ and on the minus component their negatives.
Multiplying the endpoint ratios from integration on the two components
and using the neck equation $(x-a)_+(x-a)_-=y_a$ gives
\[
 Q_{ii}\sim\frac{y_{a_0}y_{a_i}}{(a_0-a_i)^4},\qquad
 Q_{12}\sim\frac{y_{a_0}(a_1-a_2)^2}
                  {(a_0-a_1)^2(a_0-a_2)^2}.
\]
The leading units are
\[ (Q_{11}/h^3,Q_{12}/h^2,Q_{22}/h^3)_{h=0}=(-v_0^3,-v_0^2,v_0^3),\qquad v_0=3/(25\sqrt5). \]

The normalized ratios are units of $\mathcal O_L[[h]]$. Indeed the
universal graph-period construction in \cite[Thms.~4.3, 5.1--5.3]{Ich}
gives multiplicative periods as units after inversion of the three node
parameters. A unit in this localization is an edge monomial times a
power-series unit, since the edge variables are prime. The exponents of
these monomials are the products of the oriented cycle multiplicities on
each edge (equivalently by the neck integration above with the three
parameters independent). To specialize the marked deformation integrally,
take the half-edge coordinates $a_i$ on one component and $a_i+1$ on
the other. Differences of adjacent marked points are units because
$\operatorname{disc}(g)=125$; opposite coordinates differ by one. The
node smoothing ideals specialize to $(h^2),(h),(h)$ by the critical-value
calculation. The graph is trivalent so its components with their three
marked nodes have no further moduli. These choices give the formal
specialization in the universal marked chart. Theorems 5.2 and 5.3
identify the same units with the periods of the polarized Mumford
Jacobian and the exponentiated complex periods.

For completeness, the integral endomorphism $S=\pi_*r\pi^*$, with $r$
a rotation of the genus-four cover, is Rosati self-adjoint. Its pullback
equals $(r+r^{-1})\pi^*$, and $\sum_{i=0}^4r^i=0$ on the Jacobian
since the rotation quotient has genus zero. Thus $S^2+S-1=0$. The
action of $R=1+S$ on the period lattice satisfies
$R^tM_0=M_0R,\ \operatorname{tr}R=1,\ \det R=-1$. The two integral
solutions are
$\left(\begin{smallmatrix}2&1\\-1&-1\end{smallmatrix}\right)$ and its
complement $1-R$. Indeed the equations give $R_{11}=2+3m$ and a
quadratic with discriminant $-20m(m+1)$, forcing $m=0,-1$.
Applying multiplicative Rosati adjunction
\cite[Prop.~3.8]{PR} gives $Q_{11}Q_{22}=Q_{12}^3$. Set
$\rho=Q_{11}/Q_{22}$ and $\theta=Q_{11}/Q_{12}$. Then
\[
 Q_{11}=\theta^3/\rho,\qquad Q_{12}=\theta^2/\rho,\qquad
 Q_{22}=\theta^3/\rho^2 .
\]
Geometric splitting occurs exactly when $\rho$ is torsion. In fact,
real multiplication forces any elliptic quotient to give a square up
to isogeny. Otherwise the real quadratic field would map unitally into
the endomorphism algebra of a single elliptic factor, which is $\Q$ or
imaginary quadratic. For an isogeny to a square, the Hom description for
toric uniformizations \cite[\S3.7]{PR} gives nonsingular integral
character and lattice matrices. Clearing their determinants makes all
three periods powers of one Tate parameter; their valuations then force
$\rho$ to be torsion. Conversely, if $\rho$ has order $N$, the $N$th
powers of the periods belong to the period lattice of
$E_{\theta^N}^2$ with exponent matrix $M_0$. The induced torus map
has finite kernel and cokernel on lattices.

At the degeneration $h=0$ the normalized branch divisor of
\eqref{branchx} has one point on each component, at $x=-3$ on $y=g(x)$
and $x=-1/3$ on $y=-g(x)$. Subtracting the infinity divisor gives
gluing units $((a+3)/(3a+1))_{a=a_1,a_2}$ after normalization at
$a_0=1$, namely $-9\mp4\sqrt5$. These are reciprocal and not roots of
unity. If $p\ge7$, put $s_p=v_p(F_{p-(5/p)})$ as above. Each unit has
$p$-adic logarithm of valuation $s_p$ by Binet's formula, now using
$-9+4\sqrt5=-\phi_+^{-6}$. For $h\in L$ with $v_p(h)>s_p$ this proves
nontorsion of the section $B-K_D$. To see the precision, its branch
and infinity sections avoid all nodes and define a Cartier divisor
in the family over $\mathcal O_L[[h]]$ of multidegree zero. The
identity formal Picard group has two parameters restricting at $h=0$
to split torus coordinates (smoothness follows from $H^2=0$ for
curves). Multiply the section by the residue-field order minus one
to place it in the formal group. Its two coordinates at $h$ differ
from those at zero by elements of $h\mathcal O_L$; the latter both have
valuation $s_p$. The evaluated coordinates hence have minimum valuation
$s_p\ge1$. On this ball the integral formal-group logarithm is injective,
by the same degree-$n$ integration estimate used above. This comparison
concerns the class of the divisor; the period obstruction for elliptic
quotients uses a different unit.


\subsection{Computing the first period coefficients}
For residue-normalized forms $(c_i+d_ix)\,dx/g$ on the plus component
one has $c_i+d_ix=(1-a_i)(x-a_j)$, $j\ne i,\ j\in\{1,2\}$.
Write $P_i(x,h)\,dx/y$ for their normalized continuations.
The branch-point variational formula in this normalization is
\[
 \frac{d\log Q_{ij}}{dh}
 =-2\sum_{f_q(e)=0}P_i(e,h)P_j(e,h)\frac{\partial_h f_q(e)}{f'_q(e)^2}.
\]
It follows by differentiating $P_i\,dx/y$ at fixed $x$ and applying
the bilinear relation to the resulting second-kind differential, with
zero $A$-periods. At a simple branch point use $x-e=t^2$: its double
pole has leading coefficient $-P_i(e)\partial_h f_q(e)/f'_q(e)^{3/2}$
and a primitive of the other form has linear coefficient
$2P_j(e)/\sqrt{f'_q(e)}$. This yields the factor $-2$.

For the three products $1,x,x^2$ in the numerators the corresponding
raw coefficients are $3/[5(q^3-1)],0,3q/[5(q^3-1)]$.
One algebraic verification is to set $B=qf_q$, and use
\[
 B_xZ-10q(q^3-1)B_q=
 B(10-22q^3+48qx-18q^2x^2-60x^3-18qx^4),
\]
where $Z=-3q^2-2(q^3-1)x+21qx^2-3q^2x^3-12x^4-3qx^5$.
At branch points, logarithmic derivatives reduce by this identity
and by Lagrange interpolation. More precisely for the monic $f_q$,
$\sum_e H(e)/f'_q(e)$ is the coefficient of $x^5$ in $H\bmod f_q$.

The normalized continuations of the two numerator polynomials $1,x$
are, to the needed precision,
\begin{align*}
 P_{1}^{\,0}&=1+\tfrac4{25}h+\tfrac{12}{625}h^2
                  -\tfrac{544}{15625}h^3+O(h^4),\\
 P_x^{\,0}&=x(1-\tfrac4{25}h+\tfrac{112}{625}h^2
                   -\tfrac{2556}{15625}h^3)+O(h^4).
\end{align*}
Here is a reproducible residue check for normalization. Expand
$f_{1+h}=g^2+s_1h+s_2h^2+s_3h^3+\cdots$ with
$s_1=4(x-x^5),\,s_2=4x^5,\,s_3=-4x^5$.
For each numerator above form $P(1+(f-g^2)/g^2)^{-1/2}/g$ by the
binomial formula to degree three in $h$. At each nonconstant power,
partial fractions in $x$ have zero simple-pole coefficients at
$a_0,a_1,a_2$: explicitly, if the term is $D/g^{2j+1}$, the coefficient
at $a$ is
\[
 [z^{2j}]\frac{D(a+z)}
                  {(g(a+z)/z)^{2j+1}} .
\]
For $j=1,2,3$ substitution of the three displayed numerator
coefficients gives zero at $a=1$ and zero modulo $a^2+3a+1$.
This fixes the $A$-periods and provides finite checks with denominators
invertible at $p\ge7$.

Applying the variational formula to $c_iP_{1}^{\,0}+d_iP_x^{\,0}$ gives
\begin{align*}
 (\log Q_{11})'={}&3/h-3/2-9\sqrt5/50
 +(173/250+9\sqrt5/50)h\\
 &\hspace{22mm}+(-36/125-2367\sqrt5/15625)h^2+O(h^3),\\
 (\log Q_{12})'={}&2/h-1+(173/375)h-(24/125)h^2+O(h^3).
\end{align*}
The other diagonal entry follows by conjugation. In particular,
\[
 \rho=-1+(9\sqrt5/25)h+O(h^2).
\]
Its linear coefficient is a unit, and the higher coefficients are
integral by the universal-period argument. Put $T=-\rho-1$; then
$h=G(T)\in T\mathcal O_L[[T]]$ with unit linear coefficient.
The involution $q\mapsto1/q$ swaps the two quadratic nodes and
exchanges the components. In the chosen graph basis it acts by minus
the swap matrix. Thus it sends $(\rho,\theta)$ to
$(\rho^{-1},\theta/\rho)$. The series $\theta/T$ is invariant under
$T\mapsto-T/(1+T)$ and has the form
\[
 \theta=T H(U),\qquad U=T^2/(1+T),\qquad
 H(0)=-1/15,\quad H'(0)/H(0)=-14/135 .
\]
Indeed the fixed subring is $\mathcal O_L[[U]]$ by
$T^2-UT-U=0$ and the involution $T\mapsto U-T$.
The displayed logarithmic derivatives and leading units give
$\theta/T=-1/15+(14/3125)h^2+O(h^3)$ and
$U=(81/125)h^2+O(h^3)$. Conjugating $\sqrt5$ induces the same exchange
of periods and fixes $H$. This descent over the formal coefficient
field gives $H\in\Z[1/30][[U]]^\times$.

\begin{proof}[Proof of Theorem~\ref{local}]
For a prescribed multiplicative elliptic quotient, $\rho$ is torsion.
Then $-\rho$ is a nontrivial $p$-power root of unity $\zeta$, say of
order $p^r$, since $h\ne0$. Write $D=(p-1)p^{r-1}$. We have
$h=G(\zeta-1)\in L(\zeta)$, $v(h)=v(\theta)=1/D$.
Commensurability of the Tate period lattices for the target gives
$\theta^{eD}/q_E$ torsion by comparison of valuations.
If $r\ge2$, clearing this root of unity by a root of unity makes a
$p$th root of $q_E$ lie in an abelian extension of $\Q_p$.
Such a root has valuation $e/p$ with denominator $p$,
so it generates a degree-$p$ extension. This extension would be Galois
and would contain $\zeta_p$, which is impossible since $p-1\nmid p$.

For $r=1$ put $\lambda=\zeta-1$ and $a=H'(0)/H(0)$. The cyclotomic
equation gives
\[
 \lambda^{p-1}/p=-(1-\lambda/2+\lambda^2/3)
                  \pmod{\lambda^3}.
\]
After dividing $\theta^{eD}/q_E$ by its residue-field root of unity,
its linear and quadratic coefficients are
$-e/2$ and $e(1/3-a)+e(e-1)/8$. Indeed
$U=\lambda^2+O(\lambda^3)$ and the remaining unramified unit factor
differs from its Teichm\"uller lift by a multiple of $p$.
If this ratio were torsion the remaining unit would be
$(1+\lambda)^k$, since $L(\zeta_p)$ contains no $p^2$-th roots.
The linear term forces $k=-e/2\bmod p$ and the quadratic term would
be $e(e+2)/8$. The discrepancy is $-e(a+1/24)$ with
$a+1/24=-67/1080$, nonzero under the theorem's hypotheses.

For the last assertion take a primitive $\zeta$ of order $p^r$ for
each $r\ge1$ and $q=1+G(\zeta-1)$. The periods give a geometric
elliptic quotient as shown above; the two-dimensional torus power
map and the period lattice embedding with matrix $M_0$ give a finite
degree bound for each fixed $r$. The parameter lies in the finite
union of the bounded-degree quotient loci used after
Proposition~\ref{fourpt}, so it is algebraic over $\Q$, despite its
initial definition by a $p$-adic power series. The distinct valuations
give infinitely many different parameters.
\end{proof}

\section{Moving the branch class on a fixed genus-two quotient}
We describe the geometric variation in a dihedral Hurwitz fibre. Work
over an algebraically closed field of characteristic zero. Fix an
\'etale cyclic-five cover $\nu:C\to B_0$ of a genus-two curve, its
rotation $r$, and a lift $\tau$ of the hyperelliptic involution. The six
Weierstrass fibres each contribute one fixed point of $\tau$. Hence
$g(C)=6$, $g(D)=2$ for $D=C/\tau$, and the quotient $\pi:C\to D$ has
six branch points. In this section write $B_{\rm ram}=\pi_*R$,
$A=\Jac(D)$, and $\delta=B_{\rm ram}-3K_D$.
If $F=f^*\mathcal O(1)$ is the degree-five pencil on $D$,
Riemann--Hurwitz for the six branch partitions $(2,2,1)$ yields
\[
 B_{\rm ram}\sim2F-2K_D,\qquad \delta\sim2F-5K_D.
\]
Indeed the full inverse image of the six branch values consists of
twice the ramification divisor $K_D+2F$ and the six unramified points.

The integral Rosati self-adjoint endomorphism $S=\pi_*r\pi^*$
satisfies $S^2+S-1=0$. In fact $\pi^*S=(r+r^{-1})\pi^*$, and
$\nu_*\pi^*=0$ since the hyperelliptic involution on $B_0$ acts
as $-1$ on its Jacobian. Thus $(1+r+\cdots+r^4)\pi^*=0$,
which proves the identity (the kernel of $\pi^*$ is finite).
Write $T=\pi_*r^2\pi^*=-1-S$, $u=1+2S$, so $u^2=[5]$.

\begin{proposition}\label{moving}
For a fixed marked quotient $(D,S)$ in the image, the embedded
one-dimensional family through a smooth cover is birational to an
open subset of $D$. With its parameter denoted by $P$, the branch
class satisfies
\[
 \delta_P-\delta_{P_0}=-2u[P-P_0]\quad\text{in }\Jac(D).
\]
Additional level or cover markings can give finite covers of this
embedded family.
\end{proposition}
\begin{proof}
Use the embedding $(\pi,\pi r^2):C\hookrightarrow D^2$ of
\cite[Lemmas 2.4--2.6]{AP}. The curve has bidegree $(2,2)$, genus six
and self-intersection two. Let $q:D^2\to\operatorname{Sym}^2D$
and let $h:\operatorname{Sym}^2D\to A$ be the Abel map
$\{x,y\}\mapsto[x+y-K_D]$, the blowup at $0$. Factor exchange on $C$
is the reflection $\tau r^2$ and its quotient map is $\pi r$.
Thus $q(C)$ is an embedded copy of $D$ (a group-invariant closed
immersion descends to a closed immersion on these finite-group
quotients). Its strict image under $h$ is a translate of $S(D)$,
where $D$ is Abel-embedded using a geometric Weierstrass point:
the induced correspondence is $\pi_*(\pi r)^*=S$. As $S$ is an
automorphism, the image has square two. Its strict transform has
square one since $q^*q(C)=C$. It consequently passes through zero
with multiplicity one, and has the unique expression $S(D-P)$ with
$P\in D$. Uniqueness uses the trivial translation stabilizer of a
principal theta divisor.

Conversely write $G_\iota$ for the graph of the hyperelliptic
involution. The divisors
\[
 C_P=(h q)^*(S(D-P))-G_\iota
\]
form an algebraic family. Near the smooth initial member, its curves
remain smooth and connected with the same bidegree, genus and cross
correspondence $T$. To check that these smoothings are again marked
dihedral covers, let $\tau_1,\tau_2$ be the two degree-two projection
involutions. The map $A^2\to\Jac(C_P)$ by the two pullbacks has Gram
matrix $\left(\begin{smallmatrix}2&T\\T&2\end{smallmatrix}\right)$
with isogeny determinant $4-T^2=3+T$. It has four-dimensional image.
There the two involutions have respective matrices
$\left(\begin{smallmatrix}1&T\\0&-1\end{smallmatrix}\right)$
and $\left(\begin{smallmatrix}-1&0\\T&1\end{smallmatrix}\right)$.
Using $T^2+T-1=0$, their product satisfies $\Phi_5=0$ on this
image and acts trivially on its two-dimensional orthogonal complement,
where both involutions act by $-1$. The rotation thus has order five.
Its invariant differentials have dimension two; Riemann--Hurwitz then
gives an \'etale quotient of genus two. The descended reflection acts
as $-1$ on its differentials and hence is hyperelliptic. This supplies
the converse on the smooth open. Translation of $S(D)$ depends
nontrivially on $P$, giving the birational parameter assertion.

Write $\mathcal O(C_P-C_{P_0})=p_1^*M\otimes p_2^*M$ on $D^2$.
Here the two factors agree by symmetry, and the difference is
algebraically trivial. With $\lambda_\Theta(z)=t_z^*\mathcal O(\Theta)
\otimes\mathcal O(\Theta)^{-1}$, restriction to the Abel curve sends
$\lambda_\Theta$ to $-1$ on $A$: on a generic translate the difference
of restricted theta divisors is $W+\iota(P)-2W$ for the chosen
Weierstrass base $W$. Since $S$ is self-adjoint, translating $S(\Theta)$
by $-S[P-P_0]$ gives $M=-S^{-1}[P-P_0]$ in additive notation.

Let $\pi_*\mathcal O_{C_P}=\mathcal O_D\oplus L^{-1}$ with $\deg L=3$.
Adjunction gives $F=K_D+L$ since $\pi^*F=K_{C_P}$ and $\pi^*$ on
$\Pic^0 D$ is injective (apply the displayed splitting to a degree-zero
line bundle whose pullback is trivial). The surface divisor sequence
of $C_P\subset D^2$ and relative duality for $p_1$ yield
\[
 L=\det(p_{1,*}(\mathcal O(C_P)\otimes p_2^*K_D))
 \quad\text{in }\Pic D
\]
up to a constant one-dimensional vector space. Indeed
$R^1 p_{1,*}\mathcal O(-C_P)$ fits into the extension of
$H^1(D,\mathcal O_D)\otimes\mathcal O_D$ by $L^{-1}$.
The direct image has rank three since the restrictions to the
fibres have degree four. Replacing $\mathcal O(C_{P_0})$ by
$\mathcal O(C_{P_0})p_1^*M p_2^*M$ changes the determinant by
$3M+T(M)$: the first term is the projection formula; the second
is the restriction correspondence on the two horizontal curves
in a divisor presentation $M=[x-y]$, by the restriction exact
sequences. Degree-zero divisors of this form generate $A$.
Finally $(3+T)S^{-1}=1+2S$, which gives the formula for $\delta$.
\end{proof}

In particular the differences of branch classes generate all of $A$.
When a degree-two map $f:D\to E$ exists, applying $f_*$ gives a
nonconstant projection. More precisely, let $\phi=f_*$ and act on
divisors by the branch correspondence of Proposition~\ref{moving}.
On the image of $u\phi^*:E\to A$, the returned elliptic endomorphism
is $-2\phi u^2\phi^*=[-20]$; applying the normalized modified
diagonal gives $[-8]$. The formulas persist over every field where
the marked correspondence and $f$ descend. One can realize the
operations using actual divisors. A hyperelliptic fibre in the smooth
parameter open, of degree two, supplies the correspondence $2T$
by summing the two embedded curves, and hence $u=-1-2T$ integrally.
For a supplied point $Q$ the divisor $f^*([Q]-[O])$ is defined over
its field of definition. A degree-zero divisor $\xi$ can be moved
into the open using $\xi+2K_D$, a base-point-free divisor class
of degree four. Thus individual closed parameters of degree at
most four suffice when the base field is infinite. This operation
returns a multiple of a supplied point. In particular if $M/K$ is a
power-of-two extension of degree at most eight defining the maps
and $p\ne2,5$, a class in $H^1(K,E[p^n])$ is Kummer precisely when
its restricted image under $u\phi^*$ is Kummer on $A(M)$.
For the reverse implication apply the returning map and corestriction:
multiplication by $-20[M:K]$ is invertible modulo $p^n$.

There is also a useful geometric check on the field bound for
\emph{unmarked} targets. For a non-CM $E$ let $t$ be an order-three
automorphism of $E[2]$. It is Galois equivariant over an extension of
degree at most two (its centralizer in $S_3$ is $A_3$). Gluing $E$ to
itself along the graph of $t$ yields a principally polarized genus-two
Jacobian with two degree-two maps to $E$ \cite[Props.~3--4]{HLP}.
Indeed $t$ is an anti-isometry and is not induced by a geometric
automorphism of $E$. The self-adjoint matrix
$\left(\begin{smallmatrix}0&1\\1&-1\end{smallmatrix}\right)$ preserves
the graph since $t^2+t+1=0$, and descends as an integral endomorphism
satisfying $S^2+S-1=0$. This construction need not carry the marked
branch point of the explicit quintic. Conversely, in an optimal
degree-$n$ gluing of $E,E'$ with such an endomorphism, let
$\pi:E\times E'\to\Jac(D)$ be the gluing map. The matrix
$\pi^\dagger S\pi=\left(\begin{smallmatrix}a&\psi^\dagger\\
\psi&d\end{smallmatrix}\right)$ satisfies $M^2+nM-n^2=0$.
The off-diagonal map is nonzero and $a+d=-n$, giving
\[
 1\le\deg\psi=n^2-na-a^2\le5n^2/4 .
\]
There are consequently only finitely many geometric $D$ of this type
for a specified non-CM target and a fixed bound on the cover degree.

\section*{Arithmetic scope}
The branch calculation gives Jacobian points arising from
modified diagonals. The two-component locus and the twist computation
separate geometric target coverage from the arithmetic trace of those
points. A relation to a specified Selmer class or to the complex Taylor
coefficient of an elliptic $L$-function would require a further
comparison.

\begin{thebibliography}{9}
\bibitem{AP} A. Albano and G. P. Pirola, \emph{Dihedral monodromy and Xiao fibrations}.
\href{https://arxiv.org/abs/1406.6308}{arXiv:1406.6308}.
\bibitem{GS} B. Gross and C. Schoen, \emph{The modified diagonal cycle on the triple
product of a pointed curve}, Ann. Inst. Fourier 45 (1995), 649--679.
\href{https://doi.org/10.5802/aif.1469}{Article}.
\bibitem{CH} H. Darmon, M. Daub, S. Lichtenstein and V. Rotger,
\emph{Algorithms for Chow--Heegner points via iterated integrals},
\href{https://wstein.org/papers/mcom2927-chow-heegner.pdf}{author-hosted paper}.
\bibitem{MS} A. Malmendier and T. Shaska, \emph{A universal genus-two curve from
Siegel modular forms}, SIGMA 13 (2017), 089.
\href{https://sigma-journal.com/2017/089/}{Article}.
\bibitem{CFG} F. Cl\'ery, C. Faber and G. van der Geer,
\emph{Covariants of binary sextics and vector-valued Siegel modular forms of genus two}.
\href{https://arxiv.org/abs/1606.07014}{arXiv:1606.07014}.
\bibitem{LS} J. Laga and A. Shnidman, \emph{Vanishing criteria for Ceresa cycles},
Compos. Math. 161 (2025), 3017--3043.
\href{https://arxiv.org/abs/2406.03891}{arXiv:2406.03891}.
\bibitem{HZ} E. W. Howe and H. J. Zhu, \emph{On the existence of absolutely simple
abelian varieties of a given dimension over an arbitrary field},
J. Number Theory 92 (2002), 139--163.
\href{https://arxiv.org/abs/math/0002205}{arXiv:math/0002205}.
\bibitem{Kumar} A. Kumar, \emph{Hilbert modular surfaces for square discriminants and
elliptic subfields of genus 2 function fields}.
\href{https://arxiv.org/abs/1412.2849}{arXiv:1412.2849}.
\bibitem{Shaska} T. Shaska, \emph{Genus 2 fields with degree 3 elliptic subfields}.
\href{https://arxiv.org/abs/math/0109155}{arXiv:math/0109155}.
\bibitem{Ich} T. Ichikawa, \emph{The universal Mumford curve, and its abelian
differentials and periods in arithmetic formal geometry}.
\href{https://arxiv.org/abs/2010.11517}{arXiv:2010.11517}.
\bibitem{PR} M. Papikian and J. Rabinoff, \emph{Optimal quotients of Jacobians
with toric reduction and component groups}.
\href{https://arxiv.org/abs/1212.3574}{arXiv:1212.3574}.
\bibitem{HLP} E. W. Howe, F. Lepr\'evost and B. Poonen, \emph{Large torsion
subgroups of split Jacobians of curves of genus two or three}.
\href{https://math.mit.edu/~poonen/papers/large.pdf}{Author's paper}.
\end{thebibliography}
\end{document}
